% introduction.tex -- Data Availability Committee Paper (RU-V6)

\section{Introduction}

Enterprise blockchain deployments require that operational data---supply chain
records, industrial process telemetry, financial transactions---remain
confidential to the originating organization, while counterparties, auditors,
and regulators must be able to verify state transition correctness without
accessing the underlying data. Zero-knowledge validium architectures address
this tension by maintaining enterprise state off-chain and submitting succinct
validity proofs to a shared settlement layer~\cite{ethereum2024validium}. The
data availability committee (DAC) is the mechanism by which off-chain data
persistence is attested: a quorum of committee members signs a certificate
confirming that the batch data is stored and retrievable, and this certificate
is verified on-chain alongside the ZK proof.

Despite the critical role of DACs in validium security, \emph{no production
DAC provides data privacy}. StarkEx~\cite{starkware2024da}, Polygon
CDK~\cite{polygon2024cdk}, and Arbitrum Nova~\cite{arbitrum2024anytrust} all
distribute the \emph{complete} batch data to every committee member. Each
member sees every transaction, every state update, every enterprise record in
the clear. Privacy relies entirely on trust in committee members---a single
compromised or malicious member can exfiltrate the entire dataset. For
enterprise deployments where the batch data contains proprietary operational
information, this is a fundamental architectural limitation. EigenDA achieves
partial privacy through erasure coding (each node sees only one
chunk)~\cite{eigenlayer2025v2}, but chunks contain partial plaintext data and
the disperser sees the complete blob before encoding.

This paper introduces \emph{Shamir-DAC}, a data availability committee design
that replaces full data replication with Shamir's $(k,n)$-threshold secret
sharing~\cite{shamir1979share} over the BN128 scalar field. The protocol
achieves \emph{information-theoretic privacy}: fewer than $k$ committee
members learn \emph{zero} information about the batch data, regardless of
their computational power. This is strictly stronger than the computational
privacy offered by encryption-based approaches (e.g., encrypt-then-erasure-code),
as it holds against adversaries with unbounded resources. To our knowledge,
this is the first DAC design that provides information-theoretic data privacy
while maintaining sub-second attestation latency for enterprise batch sizes.

Our experimental evaluation demonstrates that Shamir-DAC achieves 175~ms P95
attestation latency for 500~KB enterprise batches in a $(2,3)$-threshold
configuration---11$\times$ below the 2-second target---with 51/51 privacy
tests confirming zero information leakage from individual shares, 61/61
recovery tests confirming fault tolerance under single-node failure, and
3.87$\times$ storage overhead (less than EigenDA's 8$\times$ expansion). All
benchmarks are implemented in JavaScript BigInt arithmetic, establishing a
conservative worst-case baseline; production deployment with native field
arithmetic is expected to improve performance by 10--50$\times$.
